\begin{beginnerstatement}[Kurz erklärt: Sicherheitsgrenzen]

\emph{Security} bezeichnet den Schutz einer Anwendung, ihrer Daten und ihrer
Zugänge vor unerlaubter Nutzung. Ein \emph{Secret} ist ein vertraulicher Wert;
\emph{Credentials} sind Zugangsdaten, mit denen sich jemand oder ein Dienst
ausweist. CORS legt fest, von welchen Web-Ursprüngen ein Browser Anfragen
zulassen darf; eine CSP begrenzt, welche Inhalte und Aktionen eine Anwendung
laden oder ausführen darf. TLS verschlüsselt die Übertragung über ein Netzwerk.
Authentifizierung beantwortet die Frage: \emph{Wer bist du?} Autorisierung
beantwortet die Frage: \emph{Was darfst du?} RBAC setzt solche Rechte über
zugewiesene Rollen um. Verifiziert sind hier nur bestimmte Konfigurations-,
Secret- und Desktop-Grenzen der Baseline, während Authentifizierung,
Autorisierung, produktive Zugangsdaten, TLS und eine vollständige
Sicherheitsbewertung Aufgabe des konkreten Produkts bleiben.

\end{beginnerstatement}
